\Needspace{13\baselineskip}
\Article{9}
In the same paper, on p. 263, the distinguished Euler tried to establish our theorem in yet another way, whose substance is as follows. Given an equation $x^n+Ax^{n-1}+Bx^{n-2}\ \text{etc.}=0$, no analytical expression representing its roots has yet been found if the exponent $n>4$; nevertheless it seems certain (as Euler asserts) that it can contain nothing but arithmetical operations and extractions of roots, combined more intricately the greater $n$ is. If this is granted, Euler shows very well that, however intricately the radical signs are combined, the value of the formula can always be represented in the form $M+N\BareRoot{-1}$, where $M,N$ are real quantities.

Against this reasoning one can object that, after the efforts of so many great mathematicians, very little hope remains of ever arriving at the general solution of algebraic equations, so that it becomes more and more probable that such a solution is altogether impossible and contradictory. This should seem all the less paradoxical \textit{since what is commonly called the solution of an equation is, strictly speaking, nothing but its reduction to pure equations}. For the solution of pure equations is not taught by this, but supposed; and if you express a root of the equation $x^m=H$ by $\BareRoot[m]{H}$, you have by no means solved it, nor have you done more than if you devised some sign to denote a root of the equation $x^n+Ax^{n-1}+\text{etc.}=0$ and put the root equal to it. It is true that pure equations far surpass all others because of the ease of finding their roots by approximation and because of the elegant connection all their roots have with one another. It is therefore by no means to be criticized that analysts have denoted their roots by a special sign. Yet from the fact that they have included this sign, just like the arithmetical signs of addition, subtraction, multiplication, division and raising to a power, under the name of \textit{analytical expressions}, it by no means follows that a root of any equation can be represented by them. Or, putting words aside, it is supposed without sufficient reason that the solution of any equation can be reduced to the solution of pure equations. Perhaps it would not be so difficult to prove this impossibility with full rigor already for the fifth degree; I shall present my investigations on this matter more fully elsewhere.
\SourceLeaf{023}
Here it is enough that the general solvability of equations, understood in that sense, is still very doubtful, and that therefore a proof whose entire force depends on that supposition carries no weight in the present state of the subject.
